\documentclass[conference]{IEEEtran}
\IEEEoverridecommandlockouts
% The preceding line is only needed to identify funding in the first footnote. If that is unneeded, please comment it out.
\usepackage{cite}
\usepackage{amsmath,amssymb,amsfonts}
\usepackage{algorithmic}
\usepackage{graphicx}
\usepackage{textcomp}
\usepackage{xcolor}
\usepackage{fancyhdr}
\usepackage{url}
\usepackage{tabularx}
\def\BibTeX{{\rm B\kern-.05em{\sc i\kern-.025em b}\kern-.08em
    T\kern-.1667em\lower.7ex\hbox{E}\kern-.125emX}}

\pagestyle{plain}
\begin{document}

\title{Ethical Discussion Between Consumer Data Regulations Across Different Jurisdictions\\
}

\author{\IEEEauthorblockN{Alvin Ton}
\IEEEauthorblockA{\textit{Western Washington University} \\
Bellingham, United States \\
tonk4@wwu.edu}
}

\maketitle
\thispagestyle{fancy}

\begin{abstract}
With the rapidly changing digital landscape in the past 20+ years, innovations
in technology spark discussion in who controls the rights of consumer data when
interacting with technology. This paper seeks to compare consumer data governance in three different jurisdictions; the U.S, EU, and China to highlight
differences in regulatory (and by extension) ethical philosophies by each country. By comparing how each jurisdiction handles consumer data, this paper
reveals the level of power that is allocated between individuals, companies, and
the state. This paper’s research reveals how the United States’s sectored, but
laxed approach to consumer data shifts the power given to companies, while the
strict and defined regulations of the European Union leads to consumer data
being in the hands of individuals. To contrast the European Union, China’s approach to consumer data is more state centric, and thus contains a philosophy
that is more government centric in how consumer data is controlled and taken.
These differences reveal many ethical questions and tensions between individual
privacy, control of data, and how the former two aspects can affect economic
growth.
\end{abstract}

\begin{IEEEkeywords}
Health Data Laws, Data Disclosure Laws, General Data Aggregation Laws
\end{IEEEkeywords}

\section{Introduction}
 Over the past 10 years, innovations in technology have completely changed the landscape of the digital world. A recurring discussed issue is how consumer data is handled. In the United States, there is a tug of war for better data transparency by disclosing how data is being stored, secured, and revealing who has access to the data for what reason. Data breaches are on the rise, and consumers are rightfully concerned with these types of questions. Individuals have pushed back by demanding how their data is being stored, what the data is being used for, and whether it is reasonable for corporations to ask for more data than what is reasonable to create the product that customers use. This paper attempts to compare laws that are currently implemented across several geo-locations to see how effective these different approaches to privacy are, and learning any tradeoffs that come with giving too much power to the state, corporations, or even to individual netizens.

Not surprisingly, there are many laws across different countries that attempt to address the same issues in their own way. The way laws are implemented solely depends on many factors, such as a country's geopolitical location and culture. This paper seeks to focus on three separate jurisdictions; the United States, the European Union, and China, and how these different regions approach data privacy and collection. The United States can be categorized as approaching problems by grouping them into de-centralized sectors, and any situation that falls within the sector is susceptible to those laws and regulations. The United States uses these sectors to dictate what corporations are allowed to do when it comes to data privacy, usage and aggregation [1]. The European Union approaches data much differently; shifting the power back to the individual consumers [1]. The EU's regulations for this reason views data and digital rights as a human right [2]. As a hybrid of both, China takes certain elements from the United States and EU by simultaneously seeking to balance companies from exploiting user data but also shifting the power to the government state itself [3]. China is best described as somewhere in the middle of the other two examples, but the difference comes from how the state is set up to request data from any individual or company at any time in the name of national security. This paper will further use examples of how each country approaches specific data, like health data, data disclosure laws, and any general data privacy rights provided to users. Each example will highlight trends in how countries approach different types of data. 

Furthermore, this paper discusses the advantages and disadvantages of each region's general approach to data. The questions that will be tackled in the end is:
\begin{itemize}
    \item What are the tradeoffs of a government that promotes a de-centralized, centralized, or a state-focused approach to data privacy and rights, and which parties end up being favored as a result of each approach?
    \item For each instance of a de-centralized, centralized, or state-focused approach, who 'loses' in each respective scenario if the proposed system is abused or misused?
\end{itemize}


\section{HEALTH DATA LAWS}


This section goes over the laws in place to address privacy in health related data. Each jurisdiction approaches health data differently. Each section will explore how each jurisdiction treats health data, and the final section will discuss the overarching between each jurisdiction and how each treat and secure health data.

The United States contains an individual sector called HIPAA (Health Insurance Portability and Accountability Act). HIPAA uses PHI (Personal Health Information) as a way to categorize health related data. HIPAA defines PHI as any and all forms of information used to identify a patient's health conditions. Furthermore PHI can be used manipulated to create, receive, or be maintained any covered entity or business associates (individuals who intend to use the data to process or treat the patient) [4]. An example of a covered entity are official healthcare providers. These providers or individuals who are working with these providers are subject under the regulations of HIPAA. Medical histories, diagnoses, and prescriptions can fall under PHI [5].

However, a current loophole in the current system are health-adjacent information that are gathered by apps or gadgets like fitness trackers, wellness applications, or menstrual apps [6]. These applications are currently not processed by covered entities, and thus are not required to comply to under HIPAA unless these applications are actively associated with covered entities or business associates [6]. If these types of data don't fall under a sector, then  the individuals data would fall under local state jurisdictions (if these jurisdictions have any laws that cover these types of data in the first place [6]). For example, Washington State happens to address health data within MHMDA (My Health My Data Act) [6]. This act treats any form of data used by covered entities or not as health data if it could be remotely used to identify or infer a patients health; either mental or physical [6]. MHMDA is further broken down into sections called RCW (Revised Code of Washington). The RCW is a standard labeling of data-related laws that are actively enforced in Washington State [7]. In the previous example, fitness trackers wouldn't fall under HIPAA regulations, but if used in the Washington State, the data would fall under MHMDA, which would then be categorized and treated under Washington's jurisdiction.

The EU defines health-related data as a subcategory of the overarching GDPR (General Data Protection Regulation). The GDPR is a standardized set of laws to address data privacy for all nations belonging to the EU [7]. Specifically when referring to health data, the GDPR categorizes much like MHMDA; where health data is any form of information personal to the patient that could relate or imply an individual's health; either mental or physical [6,7,8]. Under the GDPR, there are certain types of data that is defined as 'special category', and health data is considered part of that definition. Special category data is types of information that require more protection considerations that standard personal data [7]. 

For data laws, China had a comprehensive major reform in this respect in the form of PIPL (Personal Information Laws) in 2021 [3]. PIPL covers many types of regulations regarding different types of data, but in the context of health data, the framework doesn't specifically define what health data is [3]. In other frameworks like in the U.S where health data is defined under HIPAA, or in the EU with the GDPR, the PIPL doesn't define health data directly, but categories of data that could be argued to be under the term of 'health data' like healthcare information and biometric identifiers are individually protected without formally being defined as 'health data' [3]. The PIPL clumps up data under either personal information (the general category), or sensitive personal information (special subcategory), which covers information like healthcare and biometrics [3]. Sensitive personal information receives much more security considerations than the standard personal information category. To access sensitive personal information, an individual must have a specific purpose for accessing and processing the information, and consent must be received from the individual that the data is meant to be kept for [3].  

When comparing how health data laws are defined and implemented within each of the mentioned countries, EU's GDPR and China's PIPL are the most similar for having a comprehensive centralized approach for defining data laws in general, and an overall strong categorization for health data laws [2,3]. Prior to China's PIPL implementation in 2021, the approach of data laws was largely fragmented, much like how the U.S approaches data with the sectored data categorization approach [1,3]. EU and China solves the issue of emerging technologies that could be categorized as health data with the GDPR and PIPL in a way that the U.S's outdated HIPAA system has difficulty accommodating for. HIPAA's oversight relies on local laws for extra coverage as opposed to updating HIPAA's regulations to match the modern landscape of data aggregation through digital apps [3]. With the current state of each regulation framework, EU and Chinese health customers are the 'winners' when it comes having control over individual personal information, meanwhile the U.S's implementation of HIPAA leaves a lot to be desired for consumers. The usage of smartphone apps like fitness trackers can be used to categorize the consumer, but how that data is handled completely depends on the local state the consumer resides in [6]. In other words, if there are no local laws to accommodate the consumer, then that data can be used in however the terms of service is specified without any extra protection under HIPAA despite arguably and technically being considered health data. 

\section{Biometric Laws}
This section goes over how biometric laws are defined and enforced in each federal jurisdiction covered. Each subsection covered is dedicated to define laws and regulations of how each country defines and approaches biometric laws, with the final section discussing the differences between each region, any ethical concerns and a discussion on what parties benefit from the implementation of these laws.

Biometric data in general is defined the exact same in each country covered in this paper [3,6,10]. However, the difference is how the data is categorized and protected. In general, biometric data is any form of data that is unique through biological or behavioral characteristics like fingerprints, iris scans, DNA, facial features, and hand-eye coordination behaviors [3,6,10]. 

The United States has separate laws and regulations depending on what sector the situation falls under. Specifically for biometric data, there is overlap in multiple sectors and are enforced differently depending on the sector. For example, in Washington state, biometric data laws are applicable to any commercial use and all policies involving collecting biometric data defined under RCW section 19.375 [6]. This law requires companies that are collecting biometric data to notify, give consent, and provide a method to opt out, while also remove the service or data upon request [8]. Furthermore, companies collecting biometric data cannot openly sell or disclose biometric data to other parties without the consent for that transaction to be permitted. In cases where biometric information needs to be shared for a specific product to function, the company must still request consent from the consumer as well as contractually guarantee that the transaction of data is confidential. In other cases where biometric data can be used like in court, the usage doesn't require consent from the user as long as the evidence was gathered legally. Biometric data collection requires reasonable handling and care, and often times can only be stored for as long as needed for the purpose of the usage of the data. Once the data is no longer needed, then the data needs to be deleted in a reasonable timeline. In cases of fraud or legal gathering for court, biometric data is allowed to be kept longer than a one-time use for authentication purposes [8,9].

At the federal U.S level, when biometric data is being handled by the DHS (Department of Homeland Security), the lines drawn around biometric data is blurred. In many instances, the DHS can gather biometric data without giving consent, and have no obligations to provide an opt-out for the service or any outward rights to request for the data to be deleted [9]. These instances are typically extreme and categorized for the purpose of national security and identification during border controls [9]. However, during typical travel and non-extreme circumstances with customs and enforcement, there are options for U.S citizens to be able to opt-out of biometric data being taken for identification [9]. 

The GDPR treats biometric data, much like health data, as a special category that requires extra precaution and protections. These extra precautions are required for gathering biometric data by requiring consent directly given to consumers, robust security for storage, and only kept and used for as long as is reasonable to perform the function the service was meant to be for [10]. When the reason for usage of the biometric data is complete, the data must be deleted. Furthermore, much like health data, users have the right to know how the biometric data is used, stored, and shared, and can request to remove the data, or opt-out of the service [10].

As mentioned briefly section 2, China's PIPL covers biometric data as sensitive personal information, which much like the GDPR's special category of data requires extra care and precautions when gathering, storing, and sharing the data. Companies must obtain consent from individuals with a very specific purpose for using biometric data [11].

However, even though China categorizes biometric data as sensitive personal information, the current enforcement of the data by companies and the state has been inconsistent [11]. The rapid technological application of biometric technology has been growing faster than the regulations can currently cover. For example, there has been instances where companies are continuing to collect biometric data more than what is needed or necessary without proper enforcement, and general privacy concerns with instances where public gatherings have been collecting facial recognition for for entry without explicit consent [11]. Although PIPL was introduced to be on top of emerging applications in technologies like biometric data, the current challenge is continuing to enforce and securing data that is constantly being used in instances that weren't thought of before [11]. 

Another concern regarding the misuse of biometric data is how China's government approaches biometric collection for security purposes. The state has large authority and power over how biometric data is collected and doesn't provide partial rights in the same was as PIPL was designed to address. PIPL contains exemptions that allow the state authorities to collect and process personal information (not just including biometric data) for national security purposes. These exceptions to the state raise privacy concerns of oppressive practices that could lead to abuse of power without any processes for checks and balances. 

A consistent trend is seen between both the U.S and China during extreme instances, the government has the right to infringe on the rights of consumers where normally the law would normally protect [11]. A concern of a system where the law doesn't apply to the government can easily gravitate towards situations of escalating oppression and unjust infringes in public trust in the government residing to protect the citizens instead of oppressing [11]. In these cases, the ability for these governments to exert power in this way favors the governments leaving citizens and consumers at the mercy of the state during extreme circumstances that could potentially erode and be more common over time. 

Disregarding the potential for authoritarianism, the framework provided by EU's GDPR and China's PIPL when approaching biometric data is robust, centralized, and for the most part up-to-date with the current evolving technologies [2,3]. These implementations of laws provide citizens residing in the country to have flexible rights over how data is managed and used over time. The U.S, however, struggles with a sectored approach, and depending on the local state can struggle with having strong data protection laws regarding biometric data. Washington State's data laws was used in this example, and is one of the best general data protection in the country, but the same cannot be said in most other states. 


\section{Data Power Dynamics}

The following section goes over the distribution of power between the three covered countries. So far there has been a consistent trend amongst the U.S, the EU, and China in how power is distributed based on the implementation and enforcement of the law. In these exchanges, specific parties have more control over the data than others, and the goal is to learn what implementation is best based on the intended based interest. 

The following table shows the strength of data laws and the protection that is given to citizens against corporations and the government:

\begin{table}[h]
\centering
\caption{Comparison of Data Rights Protections}
\begin{tabularx}{\columnwidth}{lXXX}
\hline
\textbf{Category} & \textbf{EU} & \textbf{U.S} & \textbf{China} \\
\hline
Protection vs Corporations & Very strong (GDPR) & Moderate, sector-based & Strong (PIPL) \\
Protection vs Government & Strong fundamental rights & Moderate constitutional protections & Limited \\
\hline
\end{tabularx}
\label{tab:data_rights_comparison}
\end{table}

To start, the United States system is sectored and thus struggles with expanding coverage to accommodate emerging technologies. HIPAA was never meant to account for mobile applications and the aggregation of health-adjacent data outside of a healthcare format, so it's difficult to accommodate new technologies into an old system [6]. The sectored approach when outdated gives openings for corporations who're often more profit driven to take advantage of consumers through practices that benefit the company over the individuals [6]. Due to the power vacuum left by an outdated government system, the 'winners' of the United States is often times companies. Companies have flexible control over how data is aggregated, stored, and shared, and governments only intervene when company practices become too manipulative or dirty. Out of the three options, the United States is the best at giving power to organizations.

The GDPR's approach of a centralized system for data protection is clean, over-encompassing and gives the power and rights to the citizens as opposed to the companies. This system forces companies to always think about the best interest of the user because not doing so would result in sanctions, lawsuits, and large fines [2]. The drawback of robust laws also means more inconvenience for both companies and the citizens. Processes for logging in require more steps to be able to access critical information, which is better for a cybersecurity perspective, but consistently annoying as a standard customer trying to navigate a product. Out of the three options, the GDPR is the best at giving power to citizens.

China's situation prior to PIPL's implementation could best be described as a similar legal framework as the United States; with an outdated and fragmented system, but after PIPL's implementation, China's data protection laws has become much more like the GDPR. However, where China divulges from both the U.S and EU comes from the idea that even though china's data protection laws are robust, the government can and has often intervened and infringed on rights that the protection laws were meant to protect against. For the reason of national security, China's system is the best at giving power to the government. 



\section {CONCLUSION}
The best data protection laws all depends on who is the best focus. Although not entirely intentionally, the United States is set up in a way where the companies hold a majority of the data protection power. In the EU, the individuals are set up specifically to hold most of the rights and power, and China also follows suit with EU, but it's widely documented that the chinese government also exercises authoritative power at times under the guise of national security. The best implementation of data laws depends entirely on who should have the best interest. 


\section*{Acknowledgment}
In the first submission, I failed to properly cite sources, so now the current submission is correcting this mistake by citing with the references through what I already displayed with my bibliography. Currently, my bibliography isn't in the correct order, and not every source is being used at the moment. The conference page is a simplified version of my paper's introduction, sample first section, and conclusion (which I haven't completed because I'm not fully wrote the paper). 

Also, many of the references that I'm going to be using is from a paper that I've created in the past, except repurposed to be used in this paper, which is a different scope.


\begin{thebibliography}{00}
\bibitem{b1} USAGov. 2024. The U.S. and Its Government. U.S. General Services Administration. Retrieved February 19, 2026 from \url{https://www.usa.gov/about-the-us}
\bibitem{b2} GDPR-Info. n.d. General Data Protection Regulation (GDPR). Retrieved March 15, 2026 from https://gdpr-info.eu
\bibitem{b3} Zhicheng He. 2022. When data protection norms meet digital health technology: China's regulatory approaches to health data protection. Computer Law \& Security Review 47, 105758. Retrieved March 16, 2026 from \url{https://www.sciencedirect.com/science/article/pii/S0267364922001017}
\bibitem{b4} U.S. Department of Health and Human Services. 2013. HHS.gov. Retrieved February 19, 2026 from \url{https://www.hhs.gov/hipaa/for-professionals/security/laws-regulations/index.html}
\bibitem{b5} U.S. Department of Health and Human Services. 2013. The HIPAA Privacy Rule. HHS.gov. Retrieved February 19, 2026 from \url{https://www.hhs.gov/hipaa/for-professionals/privacy/index.html}
\bibitem{b6} Washington State Legislature. 2023. Chapter 19.373 RCW: Washington My Health My Data Act. Retrieved February 19, 2026 from \url{https://app.leg.wa.gov/RCW/default.aspx?cite=19.373}
\bibitem{b7} Intersoft Consulting. 2018. Art. 9 GDPR – Processing of Special Categories of Personal Data. General Data Protection Regulation (GDPR). Retrieved February 19, 2026 from \url{https://gdpr-info.eu/art-9-gdpr/}
\bibitem{b8} Alvin, T., 2025. Comparing Laws From 3 Separate Jurisdictions: Local
Washington State, Federal U.S, and the European Union, Bellingham. 
\bibitem{b9} U.S. Department of Homeland Security. 2025. Collection and Use of Biometrics by U.S. Citizenship and Immigration Services (90 FR 49062). Retrieved March 15, 2026 from \url{https://www.federalregister.gov/documents/2025/11/03/2025-19747/collection-and-use-of-biometrics-by-us-citizenship-and-immigration-services}
\bibitem{b10} GDPR-Info. n.d. Article 9 – Processing of special categories of personal data. Retrieved March 15, 2026 from \url{https://gdpr-info.eu/art-9-gdpr}
\bibitem{b11} Graham Greenleaf. 2021. China’s Personal Information Protection Law: Towards a new global standard for data privacy? Computer Law \& Security Review 6, 2 (2021), 111–123. Retrieved March 16, 2026 from \url{https://brill.com/view/journals/clsr/6/2/article-p111_001.xml}

\end{thebibliography}
\vspace{12pt}

\end{document}
